\section{Limitations and broader impact}

\sys is a bounded benchmark, not a claim that commodity hardware can substitute for production-scale evaluation. The experiments run on one local operator machine, so the results do not characterize distributed retrieval, GPU indexing, cloud variability, second-machine portability, or multi-tenant production workloads. S3 uses \hotpotmaxionbench{}, a frozen paragraph-level proxy for full-Wikipedia multi-evidence retrieval, and S2 measures top-k retrievability after inserts rather than lower-level index visibility. Engine-level quality differences are configured-search differences under the MaxionBench adapter settings, not inherent database quality. The S3 paired audit uses FAISS exact FlatIP as a local reference for the bge-base comparison, but the archived matrix does not sweep higher service-engine settings such as 128/256 probes or ef values. These choices keep the benchmark reproducible and one-day executable while narrowing the scope of the claims.

The uncertainty estimates are conditional on this fixed benchmark setting. Confidence intervals are query- or event-level summaries, not hardware-population intervals, and \maxrowp latency is not a production tail-latency guarantee. S1 and S2 change top-1 recommendations between B0 and B2, so a screening run is not enough to make a final deployment decision for those workloads; S3 keeps the same top-1 recommendation but has low full-rank correlation. The broader impact is methodological: a reproducible benchmark can help operators choose lower-cost, lower-latency retrieval infrastructure without hiding quality or post-insert retrievability tradeoffs, but readers may overread the winner as a universal recommendation. The manuscript mitigates that risk by reporting exclusions, limitations, budget stability, and objective sensitivity directly.
